% Part:sets-functions-relations
% Chapter: size-of-sets
% Section: reduction-alt

\documentclass[../../../include/open-logic-section]{subfiles}

\begin{document}

\olfileid[pt]{sfr}{siz}{red-alt}

\olsection{Redução}

\begin{editorial}
  Esta secção prova a não enumerabilidade por redução, alinhada com os resultados
  na \olref[nen-alt]{sec}. Uma versão alternativa, ligeiramente mais elaborada,
  alinhada com os resultados na \olref[nen]{sec}, encontra-se em
  \olref[red]{sec}.
\end{editorial}

Provamos que $\Bin^\omega$ é !!{nonenumerable} por um argumento de
diagonalização. Usamos um argumento semelhante para mostrar que $\Pow{\Nat}$ é
!!{nonenumerable}. Há outra forma de provar que $\Pow{\Nat}$ é
!!{nonenumerable}: mostrar que \emph{se $\Pow{\Nat}$ é !!{enumerable} então
$\Bin^\omega$ também é !!{enumerable}}. Como sabemos que $\Bin^\omega$ é
!!{nonenumerable}, segue que $\Pow{\Nat}$ também o é.

Isto chama-se \emph{reduzir} um problema a outro. Neste caso, reduzimos o problema
de enumerar $\Bin^\omega$ ao problema de enumerar $\Pow{\Nat}$. Uma solução para
este último---uma enumeração de $\Pow{\Nat}$---daria uma solução para o
primeiro---uma enumeração de $\Bin^\omega$.

Para reduzir o problema de enumerar um conjunto~$B$ ao de enumerar um
conjunto~$A$, indicamos uma forma de transformar uma enumeração de~$A$ numa
enumeração de~$B$. A forma mais simples é definir !!a{surjection} $f\colon A \to
B$. Se $x_1$, $x_2$, \dots{} enumeram~$A$, então $f(x_1)$, $f(x_2)$, \dots{}
enumerariam~$B$. No nosso caso, procuramos !!a{surjection} $f\colon \Pow{\Nat}
\to \Bin^\omega$.

\begin{prob}
Mostre que, se existir uma função !!{injective} $g\colon B \to A$, e $B$~for
!!{nonenumerable}, então $A$ também o é. Faça-o mostrando como pode usar~$g$ para
transformar uma enumeração de~$A$ numa de~$B$.
\end{prob}

\begin{proof}[Prova de {\olref[nen-alt]{thm:nonenum-pownat}} por redução]
Para a redução, suponhamos que $\Pow{\Nat}$ é !!{enumerable} e portanto que existe
uma enumeração $N_{1}$, $N_{2}$, $N_{3}$, \dots

Definimos a função $f \colon \Pow{\Nat} \to \Bin^\omega$ pondo $f(N)$ igual à
cadeia $s_{k}$ tal que $s_{k}(n) = 1$ sse $n \in N$, e $s_k(n) = 0$ caso contrário.

Isto define claramente uma função, pois sempre que $N \subseteq \Nat$, qualquer
$n \in \Nat$ ou é !!a{element} de $N$ ou não é. Por exemplo, o conjunto
$2\Nat = \Setabs{2n}{n \in \Nat} = \{0,2, 4, 6, \dots\}$ dos naturais pares é
mapeado para a cadeia $1010101\dots$; $\emptyset$ é mapeado para $0000\dots$;
$\Nat$ é mapeado para $1111\dots$.

Também é !!{surjective}: toda a cadeia de $0$s e~$1$s corresponde a algum
conjunto de números naturais, nomeadamente aquele que tem como membros os naturais
nas posições em que a cadeia tem um~$1$. Com mais precisão, se $s \in \Bin^\omega$,
defina-se $N \subseteq \Nat$ por:
\[
N = \Setabs{n \in \Nat}{s(n) = 1}
\]
Então $f(N) = s$, como se verifica consultando a definição de~$f$.

Consideremos agora a lista
\[
f(N_1), f(N_2), f(N_3), \dots
\]
Como $f$ é !!{surjective}, todo o membro de $\Bin^\omega$ tem de aparecer como
valor de~$f$ para algum argumento, e portanto na lista. Esta lista teria de
enumerar todo~$\Bin^\omega$.

Logo, se $\Pow{\Nat}$ fosse !!{enumerable}, $\Bin^\omega$ seria !!{enumerable}.
Mas $\Bin^\omega$ é !!{nonenumerable} (\olref[nen-alt]{thm:nonenum-bin-omega}).
Logo $\Pow{\Nat}$ é !!{nonenumerable}.
\end{proof}

%\begin{explain}
%It is easy to be confused about the direction the reduction goes in.
%For instance, !!a{surjective} function $g \colon \Bin^\omega \to X$
%does \emph{not} establish that $X$ is !!{nonenumerable}.  (Consider $g
%\colon \Bin^\omega \to \Bin$ defined by $g(s) = s(1)$, the function
%that maps a sequence of $0$'s and $1$'s to its first !!{element}.  It
%is surjective, because some sequences start with $0$ and some start
%with $1$. But $\Bin$ is finite.)  Note also that the function $f$ must
%be surjective, or otherwise the argument does not go through:
%$f(x_1)$, $f(x_2)$, \dots{} would then not be guaranteed to include
%all the !!{element}s of~$Y$. For instance, $h\colon \Nat \to
%\Bin^\omega$ defined by
%\[
%h(n) = \underbrace{000\dots0}_{\text{$n$ $0$'s}}
%\]
%is a function, but $\Nat$ is !!{enumerable}.
%\end{explain}

\begin{prob}\label{sfr:siz:red:prob:nat-nat}
  Mostre que o conjunto~$X$ de todas as funções $f\colon \Nat \to \Nat$ é
  !!{nonenumerable} por um argumento de redução (Sugestão: dê uma função
  sobrejectiva de $X$ para~$\Bin^\omega$.)
\end{prob}

\begin{prob}
Mostre que o conjunto de todos os \emph{conjuntos de} pares de números naturais,
O.s.\@, $\Pow{\Nat \times \Nat}$, é !!{nonenumerable} por um argumento de redução.
\end{prob}

\begin{prob}
Mostre que $\Nat^\omega$, o conjunto das sequências infinitas de números
naturais, é !!{nonenumerable} por um argumento de redução.
\end{prob}

%\begin{prob}
%Let $P$ be the set of functions from $\Nat$ to the set $\{0\}$, and let $Q$ be the set of \emph{partial}
%functions from the set of positive integers to the set $\{0\}$. Show
%that $P$~is !!{enumerable} and $Q$~is not. (Hint: reduce the problem
%of enumerating $\Bin^\omega$ to enumerating~$Q$).
%\end{prob}

\begin{prob}
Seja $S$ o conjunto de todas as !!{surjection}s de $\Nat$ para o conjunto
$\{0,1\}$, O.s.\@, $S$ consiste em todas as !!{surjection}s~$f \colon \Nat \to
\Bin$. Mostre que $S$ é !!{nonenumerable}.
\end{prob}

\begin{prob}
Mostre que o conjunto~$\Real$ de todos os números reais é !!{nonenumerable}.
\end{prob}

\end{document}
